\begin{pastexam}[label={px:ntu-112-ee-ds-12}]{臺大電機所 112 學年度 第 12 題}
（複選；每題 15 分，題內每個選項單獨計分，答對得 3 分，答錯倒扣 3 分，倒扣至複選題總分 0 分為止；未作答得 0 分，不倒扣）

Which of the following statement(s) about the \textit{tree} data structure is (are) correct?
\begin{enumerate}[label=(\Alph*)]
\item (Number of edges) = (Number of nodes) - 1
\item (Number of leaf nodes) = (Number of internal nodes) + 1\\
// Note: Leaf nodes are those with no child node. Nodes that are not leaf nodes are internal nodes.
\item (Number of paths from the root to the leaf nodes) is O(Number of leaf nodes).
\item (Height of a \textit{tree}) is $\Omega$(log(Number of nodes)).
\item (Number of subtrees) = O((Number of nodes)$^2$)\\
// A \textit{subtree} \texttt{T}$_{\texttt{s}}$ of a tree \texttt{T} is a \textit{tree} formed by a subset of nodes in \texttt{T} and with the edges that connect these nodes.
\end{enumerate}
\end{pastexam}

\begin{pastsoln}
答案是 (A)(C)。

以下的「樹」是一般的有根樹（子節點個數不限），$n$ 是節點數。
\begin{itemize}
\item (A) \textbf{對}。 除了根，每個節點恰好有一條邊連到它的父節點，邊和非根節點一一對應，邊數 $= n - 1$
（$n \ge 1$；資結書第~7.1~節）。
\item (B) \textbf{錯}。 只有\textbf{滿}二元樹（每個內部節點恰好兩個子節點）才有「葉 $=$ 內部 $+ 1$」（資結書定理~7.2：
$n_0 = n_2 + 1$）。 反例：三個節點的一條鏈，葉節點 $1$ 個、內部節點 $2$ 個；或一個根帶三個子節點，葉 $3$ 個、
內部 $1$ 個。
\item (C) \textbf{對}。 從根到某一片葉子的路徑只有一條（樹沒有環），所以路徑數剛好等於葉節點數，當然是
$O(\text{葉節點數})$。
\item (D) \textbf{錯}。 一般的樹可以很矮：一個根帶 $n - 1$ 個子節點（星形），高度是 $1$，不是
$\Omega(\log n)$。 （若限定是二元樹，第 $i$ 層最多 $2^i$ 個節點，高度 $\ge \lfloor \lg n \rfloor$，就會成立；但題目
說的是 tree。）
\item (E) \textbf{錯}。 照題目的定義，子樹是「一部分節點和連接它們的邊形成的樹」，也就是一組彼此連通的節點。
星形樹中，根加上任意一部分葉節點都是連通的，光這些就有 $2^{n-1}$ 個子樹，是指數個，不是 $O(n^2)$。
\end{itemize}
\end{pastsoln}
